% Version: 2026.1.0.0
% WHAT: A formal companion to Engineering Studio AI's OpenAI Dev Week paper.
% WHY: A failed in-place write can erase an otherwise recoverable run snapshot.
% HOW: Model one namespace entry and prove atomic visibility by trace induction.
\documentclass[12pt]{article}
\usepackage[utf8]{inputenc}
\usepackage[T1]{fontenc}
\usepackage{times}
\usepackage[margin=1in]{geometry}
\usepackage{amsmath,amssymb,amsthm}
\usepackage{booktabs,tabularx,ragged2e}
\usepackage{enumitem}
\usepackage[final]{microtype}
\usepackage{fancyhdr,lastpage}
\usepackage[hidelinks]{hyperref}
\hypersetup{
  pdftitle={Atomic Publication of Engineering Studio Run Snapshots},
  pdfsubject={Engineering Studio AI formal companion}
}
\pagestyle{fancy}
\fancyhf{}
\lhead{Engineering Studio AI}
\rhead{Atomic run snapshots}
\cfoot{\thepage\ of \pageref{LastPage}}
\setlength{\headheight}{15pt}
\setlength{\emergencystretch}{2em}
\setlist[enumerate]{leftmargin=*}
\newtheorem{theorem}{Theorem}
\newtheorem{definition}[theorem]{Definition}
\newtheorem{counterexample}[theorem]{Counterexample}
\newcommand{\Bytes}{\mathcal{B}}
\newcommand{\Valid}{\mathcal{J}}
\newcommand{\absent}{\bot}

\title{Atomic Publication of Engineering Studio Run Snapshots\\
\large A formal companion to the OpenAI Dev Week paper}
\date{2026-11-20}

\begin{document}
\maketitle

\begin{abstract}
Engineering Studio restores run history from a JSON sidecar. The baseline
writer opens that sidecar directly, so serializing writes with a process-local
lock still permits a failed or interrupted write to erase a previously
complete snapshot. This companion characterizes a narrow alternative:
serialize the snapshot, write and close a private temporary file in the
destination directory, and publish it by atomic replacement. An inductive
trace model proves that a successful lookup selects the last whole
publication or observes absence before the first publication. A
counterexample shows why direct truncation lacks this property. The proof
treats filesystem semantics as explicit premises and separates file
visibility from power-loss durability, validation of arbitrary JSON,
multi-process consistency, artifact state, and API events.
\end{abstract}

\section*{Key Words}
Atomic replacement; run history; linearization point; failed publication;
snapshot isolation; process restart.

\section*{Executive Summary}
The target is one run's \texttt{run.json}. Prepare complete bytes away
from the name used by recovery. A successful rename changes which snapshot
that name selects; failures before that operation leave the entry and its
old bytes alone. Recovery addresses the final name exclusively. The
\texttt{RunStore} tests should exercise existing and absent targets, inject
write and replacement failures, and verify that old snapshots remain
loadable. The implementation maintainer owns the filesystem premise;
an independent reviewer checks the actual write, close, replace, and
cleanup order against the model. A local check on one filesystem gives
evidence for that environment.

\section*{Key Findings}
Publication is an event in the filesystem namespace, whereas a successful
return is an observation by the caller. An interrupted caller may miss the
result of an already completed publication. A cleanup failure can leave
private bytes behind without changing which bytes a recovery lookup selects.

\begin{table}[!htbp]
\centering\small
\caption{Atomic run snapshot findings, Part I}
\label{tab:findings-1}
\begin{tabularx}{\linewidth}{@{}p{0.21\linewidth}X l@{}}
\toprule
Case & Result and limit & Source \\
\midrule
Typed contents & A finite byte string differs from a missing entry, even when the string is empty. & \ref{def:states} \\
Trace cut & A lookup selects an entry at one transition boundary under the reader premise. & \ref{def:trace} \\
Zero publications & Recovery sees the initial snapshot, which can be absent. & \ref{thm:visibility} \\
First success & A successful replacement can publish to a previously missing name. & \ref{thm:visibility} \\
Write or close error & A prepublication error preserves the prior entry under staging isolation. & \ref{thm:visibility} \\
Replace error & Preservation requires failure before the namespace change. & \ref{thm:visibility} \\
\bottomrule
\end{tabularx}
\end{table}
\begin{table}[!htbp]
\centering\small
\caption{Atomic run snapshot findings, Part II}
\label{tab:findings-2}
\begin{tabularx}{\linewidth}{@{}p{0.21\linewidth}X l@{}}
\toprule
Case & Result and limit & Source \\
\midrule
Equal bytes & Two successful publications may be indistinguishable by content. & \ref{thm:visibility} \\
Competing reads & An open before replacement can finish afterward with the older bytes. & \ref{thm:visibility} \\
Cleanup error & Orphans remain invisible to the literal-name lookup. & \ref{thm:visibility} \\
Process interruption & Visibility follows the completed namespace event, not caller acknowledgement. & \ref{thm:visibility} \\
Serial writes & The last successful replacement wins for one name; no logical conflict protocol follows. & \ref{thm:visibility} \\
Direct truncation & A reader can select empty or partial bytes between open and close. & \ref{cex:truncate} \\
\bottomrule
\end{tabularx}
\end{table}

\begin{table}[!htbp]
\centering\small
\caption{Symbols used in the trace and visibility claim}
\label{tab:symbols}
\begin{tabularx}{\linewidth}{@{}lX l@{}}
\toprule
Symbol & Domain and meaning & Source \\
\midrule
$\Bytes$ & Finite byte strings, including the empty string $\epsilon$. & \ref{def:states} \\
$\Valid$ & Subset of $\Bytes$ with the required full JSON serialization. & \ref{def:states} \\
$\absent$ & Absence of the public namespace entry. & \ref{def:states} \\
$T$ & Finite trace length, measured in transitions. & \ref{def:trace} \\
$k$ & Cut index in $\{0,\ldots,T\}$. & \ref{def:trace} \\
$P_k$ & Public entry contents evaluated at cut $k$. & \ref{def:trace} \\
$S_k$ & Finite private staging map at cut $k$. & \ref{def:trace} \\
$L_k$ & Payload of the most recent commit, or the initial public value. & \ref{def:trace} \\
$A_k$ & Count of successful publications through cut $k$. & \ref{def:trace} \\
$u$ & A unique temporary namespace key. & \ref{def:states} \\
$b$ & Complete serialized attempt payload in $\Valid$. & \ref{def:states} \\
$r$ & A reader's successful lookup cut. & \ref{def:trace} \\
\bottomrule
\end{tabularx}
\end{table}
\clearpage
\tableofcontents
\listoftables
\clearpage

\section{Source and scope}
\label{sec:source}
The OpenAI Dev Week formal paper is the companion baseline for the
workflow and its lifecycle. The concrete persistence entry point is
\texttt{RunStore.\allowbreak\_persist\_locked}; the recovery projection is
\texttt{RunStore.\allowbreak load\_from\_disk} in the runtime
\texttt{runs.py} module. The baseline implementation at
\texttt{535f5de} writes directly to \texttt{run.json}. The process lock
prevents interleaving among cooperating callers within that process,
but an interruption can still occur after truncation.

This companion fixes the semantic target for a replacement writer:
keep the same public file name, serialized schema and best-effort error
behavior. It covers one run directory, without changing the pipeline,
artifact storage, history schema or API and stream payloads. Atomic
visibility is a claim about which complete bytes a directory entry can
select. The Python replacement API and POSIX rename semantics motivate
the premises below \cite{python, posix}. A deployed filesystem must
satisfy those premises; calling the API is not itself evidence about
every configured storage backend.

\section{Model and main result}
\label{sec:model}
\begin{definition}[Namespace and eligible payloads]
\label{def:states}
Let $\Bytes=\{0,\ldots,255\}^{<\omega}$ be the set of finite byte
strings, with empty string $\epsilon$. Let $\Valid\subseteq\Bytes$
be the complete UTF-8 JSON encodings of valid application snapshots.
Write $\absent\notin\Bytes$ for a missing entry. The final name is
the literal \texttt{run.json} entry of a fixed directory. Private names
$u$ are distinct from the final name. A stage is a private entry with
a byte string and a phase, open or closed. A closed stage may be
published only after its bytes have been established equal to the
attempt's fixed payload $b\in\Valid$.
\end{definition}

\begin{definition}[Trace and observation cut]
\label{def:trace}
A finite trace has $T\geq0$ transitions, indexed $1,\ldots,T$.
Cut $k$ is immediately after transition $k$, with $0$ the initial cut.
At cut $k$, $P_k\in\Bytes\cup\{\absent\}$ is the public entry value
and $S_k$ is the finite map of private entries. A publication is the
namespace transition consuming a closed stage and changing the final
entry to its complete payload. Let $A_k$ count such transitions through
cut $k$. Let $L_k$ be their last payload, or $P_0$ if $A_k=0$.
A successful reader selects the public entry at one cut $r$ within
its lookup interval, and reads the selected immutable file to the end,
or observes $\absent$ at that cut.
\end{definition}

\begin{theorem}[Atomic visibility of one run snapshot]
\label{thm:visibility}
Under the hypotheses below, every cut satisfies
\begin{equation}
  P_k=L_k\in\Valid\cup\{\absent\}, \qquad 0\le k\le T .
\end{equation}
A successful lookup observes $L_r$.
An attempt which ends before its publication changes no public bytes.
Its cleanup can affect only its private name. When an attempt
publishes, the next cut has its complete payload, regardless of
whether the caller later receives an acknowledgement.
\end{theorem}

\noindent\textbf{Proof type:} Direct, using induction on finite traces
and case analysis of the observable transitions.

\noindent\textbf{Hypotheses.}
\begin{enumerate}[label=H\arabic*.,ref=\arabic*]
\item\label{hyp:init} The initial final entry is absent or holds an
eligible payload. No external writer mutates the final path or its
file object. The directory and final name stay fixed.
\item\label{hyp:private} Cooperating writers for the fixed run are
serialized. Each attempt gets a fresh exclusive private name on the
same filesystem and never modifies another stage. Bytes are fixed
before staging. Only exact, fully written, closed stages may publish.
\item\label{hyp:replace} Successful replacement changes the public
namespace in one transition and consumes its stage name. An error
classified as prepublication leaves the public entry untouched. Cleanup
targets only the attempt's private name.
\item\label{hyp:read} Lookup follows one literal final name and selects
an absent entry or one file object atomically. A successful read returns
all bytes of that object, whose content remains immutable across later
replacements. I/O errors are distinct from observations of absence.
\item\label{hyp:trace} The trace is finite and includes the listed
preparation, publication, lookup, cleanup, interruption, and local
bookkeeping events. Process interruption can stop an attempt and leave
its private files behind. Filesystem corruption and power failure are
outside the model.
\end{enumerate}

\noindent\textbf{Proof dependencies.}
\begin{enumerate}
\item Definitions~\ref{def:states} and~\ref{def:trace} fix
eligible byte strings, event indices, and the observation cut.
\item Induction over $\{0,\ldots,T\}$ and exhaustive case splits
prove the trace invariant from H\ref{hyp:init}--H\ref{hyp:trace}.
\item H\ref{hyp:read} identifies completed reader output from its
lookup cut; H\ref{hyp:replace} supplies the publication point.
\end{enumerate}

\begin{proof}
We induct on the cut index $k$ of the fixed finite trace.
\begin{enumerate}
\item \textbf{Domains, initial condition, and candidate invariant.}
Fix $T\in\mathbb{N}_0$ and a trace satisfying H\ref{hyp:init}
through H\ref{hyp:trace}. The trace may have length zero.
A successful publication must come from an attempt that has fixed
a finite payload $b$ in $\Valid$, by H\ref{hyp:private}.
The payload is fixed before the first stage write, so a private
prefix and a payload are distinct state concepts even if their
current byte values happen to be equal.
The former can still be changed; the latter is the target of
the serialization operation.
The variable $P_k$ is a public file value, whereas $S_k$
represents entries that the reader's address does not select.
A failed lookup is modeled separately from the value $\absent$.
For example, a permissions error supplies no content observation.
If a successful lookup returns missing, it supplies an observation
of $\absent$ under H\ref{hyp:read}.
For the induction define the set
\begin{equation}
\label{eq:commits}
C_k=\{i:1\leq i\leq k,\ i\text{ is a successful publication}\}.
\end{equation}
Its finiteness follows from that of the index set.
Its cardinality is $A_k$ by the counting definition.
If $C_k$ is nonempty, the usual order of integers gives a
unique maximum, so the last payload $L_k$ is well defined.
If $C_k$ is empty, $L_k$ has the initial value by definition.
For a private key, equality of file values does not imply
equality of keys. In particular and independently of values,
H\ref{hyp:private} excludes the public key from the domain of
$S_k$ at every cut.
We will carry the invariant
\begin{equation}
\label{eq:inv}
  P_k=L_k\in\Valid\cup\{\absent\},
  \qquad
  0\leq A_k\leq k .
\end{equation}
The inequality will make the endpoints of the count explicit.
No inequality about timestamps, payload sizes, or application
status is part of this invariant.

\item \textbf{Base case.}
At $k=0$ there is no integer satisfying $1\leq i\leq k$.
Thus $C_0$ is empty, which yields $A_0=0$ and $L_0=P_0$.
H\ref{hyp:init} gives the membership part of \eqref{eq:inv}.
A prior private orphan, if present, has a key different from the
public name. Freshness includes distinction from such live keys.
A model state with an absent public entry and a nonempty private
map is therefore permitted.
In that state an initial lookup returns absence, even when
a closed private file already contains a complete JSON document.
If $T=0$, this base case supplies the entire trace conclusion
and no existence of a writer operation is inferred.
If $T>0$, it supplies the hypothesis for the first transition.
The bound is tight at this initial cut because $0\leq0\leq0$.
The induction does not require the initial byte string to have
been produced by this implementation; it requires only the
eligibility specified in H\ref{hyp:init}.

\item \textbf{Induction hypothesis and partition of the next step.}
Fix $k<T$ and assume \eqref{eq:inv} at cut $k$.
The successor index $k+1$ exists and is at most $T$.
By H\ref{hyp:trace} the next event is in one of three classes.
The first changes the public entry by a successful publication.
The second changes only preparation, staging, or attempt bookkeeping.
The third is a reader action.
If an implementation reports a successful operation after the
replacement, model the replacement and the report as distinct
events. Thus a report belongs to the second class.
If an interruption occurs immediately after replacement,
there is still a completed publication in the first class before
the interruption in the second class.
Under the atomic namespace premise an interruption cannot split a
single publication transition into two observable public updates.
A reader which starts before that transition and finishes after
it still has a selection cut at either side of it.
Distinguish event classification by the operation performed,
before considering whether byte values changed.
A replacement with equal bytes belongs to the first class,
which prevents overlap with a nonpublication that also leaves
the observed byte value equal.
One event cannot both commit a private stage and be an ordinary
read event, since the model decomposes execution at the
namespace selection and replacement boundaries.
This yields disjoint cases and covers the allowed events.
An unmodeled external unlink would require another case and
is excluded by H\ref{hyp:init}.

\item \textbf{Preparation events.}
First consider a computation of the whole serialized payload.
Its write set is local memory; there is no operation on $P_k$.
If serialization raises, no publishable stage has been
established. The state at the next cut therefore has
$P_{k+1}=P_k$ and $C_{k+1}=C_k$.
If creation of the private directory or file fails instead, it
cannot obtain the public name, by H\ref{hyp:private}.
It may have no private effect or leave some private allocation.
Either is compatible with the invariant, which depends only on
the public entry and on successful publications.
If creation succeeds, freshness adds a private key distinct
from all current keys. It does not rebind the final name.
A write of a proper prefix updates that private entry.
A write of all bytes also updates only that private entry.
Neither of these writes is a publication.
This distinction matters for a failure during buffering, which
can leave a private prefix not recognizable as JSON.
On a short-write failure the attempt cannot use that entry
as a closed stage with proved equality to $b$.
A successful close may establish exactly that equality and the
closed phase. It does not transfer the private name.
A close failure has the same public effect as a failed write,
even if all bytes in fact reached system buffers.
Its uncertain terminal stage is never offered for publication.
This relies on control ordering: reaching the replacement
requires a completed successful write and close.
A process may stop at any of these finer boundaries.
In all such preparation cases the history $C_k$ is unchanged.
Hence $A_{k+1}=A_k$ and $L_{k+1}=L_k$.
Substitution of $P_k=L_k$ proves the equality at $k+1$,
and membership follows from the induction hypothesis.
Finally $A_{k+1}=A_k\leq k<k+1$ proves the upper count bound.

\item \textbf{Replacement success and exact bytes.}
Now take the first event class. H\ref{hyp:private} supplies
one private name $u$ in the closed phase whose bytes equal
the whole payload $b\in\Valid$.
By H\ref{hyp:replace}, the replacement has one namespace effect:
it sets the final entry to that closed file and consumes the
private file name. A surrounding private directory may survive
until its own cleanup; that directory is not the public entry.
There is no cut with an intermediate public prefix.
If the old entry existed, the new file replaces that binding.
If it was absent, the new file creates the public binding.
These are the two entry states by Definition~\ref{def:states}.
At this successor cut,
\begin{align}
\label{eq:publish}
     P_{k+1} &= b,\\
     C_{k+1} &= C_k\cup\{k+1\}.
\end{align}
Because $C_k\subseteq\{1,\ldots,k\}$, $k+1$ was not already
counted. It is larger than each element of $C_k$ and is
the unique maximum of $C_{k+1}$.
Consequently $A_{k+1}=A_k+1$ and $L_{k+1}=b$.
Equation~\eqref{eq:publish} yields $P_{k+1}=L_{k+1}$.
The membership claim follows from $b\in\Valid$.
The count bound follows by adding one to $A_k\leq k$,
which also preserves nonnegativity.
A replacement of bytes $b$ by $b$ still adds precisely one
index to the commit set. A trace count can increase while the
bytes and the byte length remain unchanged.
A reader cannot infer an execution count from byte equality.

\item \textbf{Failed replacement, cleanup, and acknowledgement.}
Consider a replacement that fails before the namespace change.
H\ref{hyp:replace} gives $P_{k+1}=P_k$.
The event index is not a successful publication;
therefore $C_{k+1}=C_k$.
The same conclusion holds for cleanup of an attempt that never
published, because cleanup addresses only that attempt's fresh
private namespace. Deleting a private entry changes $S_k$ without
changing the literal final name.
A cleanup error may leave $S_k$ larger than it would otherwise
have been. Its footprint is still private.
If cleanup runs after replacement, it cannot delete the public
file by removing the original private name:
that name was consumed by replacement, and the cleanup targets
only the private temporary location.
A private directory and a public file may refer to the same
filesystem but their namespace addresses remain distinct.
The same reasoning covers later logging or return to the caller.
These operations can change warning text or caller state;
they add no index to $C_k$.
If a caller encounters an error after publication, then the
attempt is not an instance of the theorem's prepublication clause.
Its replacement still occurs in the history.
In particular it would be incorrect to infer rollback from
a warning emitted while cleaning a temporary directory.
For all these second-class events, $L_{k+1}=L_k$ and the
bounds and membership in \eqref{eq:inv} transfer unchanged.
A filesystem whose error return can report failure after an
unknown publication cannot use that return alone to establish
the failed-before-change premise.

\item \textbf{Reader events and induction closure.}
A successful read neither publishes nor writes the private map.
An unsuccessful read also cannot change either namespace
according to H\ref{hyp:read}.
Thus $P_{k+1}=P_k$ and $C_{k+1}=C_k$ for reader events.
A result from the reader may be returned at a later cut.
That local return is bookkeeping, still without a publication.
The induction hypothesis gives $P_{k+1}=L_{k+1}$ for
this third event class, with unchanged membership and bounds.
We have proved the induction step for every possible transition.
By finite induction, \eqref{eq:inv} holds for all cuts,
including the maximal index $T$.
This proves a prefix property: applying the same argument
to a trace prefix gives exactly the same $P_k$ and $L_k$.
Future transitions do not retroactively change a prior cut.
The prefix property handles a stopped execution without
requiring any future retry or eventual publication.
It also permits a read-only trace to have many events but
no commit, since none of the reader cases increments the count.


\item \textbf{From namespace state to the reader's result.}
Let a particular successful reader have lookup cut $r$.
H\ref{hyp:read} supplies $0\leq r\leq T$ and identifies
the selected entry with $P_r$. If $P_r=\absent$, there is
no selected byte string; the observation is absence.
Otherwise the reader gets a file object with contents $P_r$.
Since publishing replaces the entry and does not write into
a published object, its read-to-end observes the same finite
string even if the read returns at a later boundary.
Therefore the result is $P_r=L_r$ in both cases.
A reader exception is outside this statement because it has
no completed contents to equate with $L_r$.
A reader need not reconstruct the namespace event order.
In the proof the cut is a witness supplied by selection semantics,
not a new timestamp field or an extra disk read.
The recovery loader's test for file existence followed by open
need not be atomic as a pair.
If it skips because a literal check sees absence, that lookup
is the absence witness. With an existing entry it may open a
newer file by the time the open occurs.
Only the selection that supplied the observed result matters.
An external deletion between checks could introduce behavior
outside H\ref{hyp:init}; a later replacement alone does not
remove the previously existing final name.

\item \textbf{Before, across, and after one replacement.}
For a single successful publication at transition $h$,
the adjacent cuts are $h-1$ and $h$, where $1\leq h\leq T$.
The invariant gives
\begin{equation}
\label{eq:adjacent}
 (P_{h-1},P_h)=(L_{h-1},b).
\end{equation}
If a read selects before $h$, it cannot obtain its bytes
from the publication at $h$. Selection at $h-1$ gives
$L_{h-1}$ exactly, including possible absence.
Selection at $h$ gives $b$ exactly.
Selection later gives $b$ until the next publication,
and the last payload thereafter.
These possibilities are exhaustive for a specific replacement:
the selected cut is either strictly below $h$, exactly $h$,
or strictly above it.
A read begun before $h$ can be delayed without selecting
until after $h$, in which case it may return the new payload.
If it selected an old object before $h$ and returns after $h$,
its result remains the old payload.
The proof assigns the atomic observation to selection,
not to invocation or return.
A reader entirely between two consecutive publications
has no alternative public value under the invariant.
A reader entirely before the first publication has only
$P_0$ as an allowed successful observation.
Multiple readers can choose different cuts and return in an
order different from their selections.
Their distinct outcomes are consistent with one ordered
sequence of namespace publications.
Sorting readers by return time would not justify a monotonic
history without an additional constraint on the reads.

\item \textbf{Exactly which attempts can change the public entry.}
Consider an attempt from selection of $b$ to its last event.
Writers for this run are serialized under H\ref{hyp:private},
so no other attempt has a publication strictly within that
interval. Reader events can overlap.
If the attempt never commits, none of its own events enters
$C_k$. All prefixes of this attempt preserve the public value
at its start. The invariant independently supplies this
fact at every cut inside the attempt.
If it commits once, preparation cuts retain that prior value.
The publication cut selects precisely its closed bytes.
Later bookkeeping and cleanup in the attempt retain those
new bytes. A second replacement is not a continuation of the
same publication: the old stage key was consumed, and a new
eligible stage needs a fresh attempt.
A retry after prepublication failure can compute a new payload
and stage it under a fresh name. Such a retry starts at a cut
whose public value is still the last committed snapshot.
The proof of the first attempt does not imply that a retry
will be attempted or succeed.
A retry after an uncertain acknowledgement can publish again
even when the first attempt succeeded. The second commit
does not falsify the invariant, because its position is in $C_k$.
An idempotency record would be required to turn that fact into
an exactly-once attempt protocol. There is no such record
among the present state variables.
Best-effort persistence is compatible with this distinction:
the in-memory run can advance after an error, while recovery
can still return the prior complete sidecar.

\item \textbf{Interruption and the namespace cut.}
Take the prefix ending at the point immediately before the
process ceases to execute.
If interruption precedes the allocation of the private namespace,
that attempt has no stage file to leave behind.
If interruption occurs while the file is open, $S_k$ can
include arbitrary partially produced bytes at the private key.
If it occurs after close and before replacement, $S_k$ can
include a complete eligible payload at that key.
Neither private condition changes the observed public entry,
because H\ref{hyp:read} selects only the final name.
If interruption occurs after replacement and before cleanup,
the public entry has the new bytes, and an empty private
directory may still be discoverable by enumeration.
If interruption occurs after cleanup and before acknowledgement,
even that private directory may be absent while the new bytes
remain public.
Thus two trace prefixes with identical caller observations
may differ in their committed state.
Conversely two prefixes with different cleanup outcomes may
have identical public states.
Both statements follow by taking the indicated cuts and
substituting their respective $C_k$ into \eqref{eq:inv}.
A fresh recovery process can encounter a private orphan from
the earlier process. This changes neither its lookup address
nor the identity of the last namespace commit.
It is a separate engineering question whether such orphan
directories are reclaimed later.
Removing them globally while another process is preparing a
stage would need an ownership rule beyond this proof.
Under H\ref{hyp:private}, all in-scope cleanup names belong
to the attempt which created them.

\item \textbf{Boundary audit for contents and extent.}
The string $\epsilon$ has length zero and is an element
of $\Bytes$ by its construction as a finite string.
The symbol $\absent$ is outside $\Bytes$ and therefore distinct
from $\epsilon$ even at the abstraction boundary.
The namespace induction itself uses only closure of the
admitted content type under copying, so a variant with
$\Valid=\Bytes$ could publish $\epsilon$ as a present file.
That variant would say nothing about JSON decoding.
In the service specialization, $\epsilon\notin\Valid$:
there is no complete JSON value in the empty text.
An empty stage file before bytes arrive is thus legitimate
private state but cannot become a successful application snapshot.
When a field of a snapshot is itself the empty string, the
whole JSON serialization can still be nonempty and eligible.
This separates the edge case of an empty application value
from that of an empty byte sequence.
There is no maximum permitted byte length in the abstract set
$\Bytes$. Every individual payload has finite length, so
each prefix in the preparation case is well defined.
A concrete allocation or capacity limit can prevent reaching
the successful stage condition; it does not produce a smaller
committed version merely by truncating private bytes.
The abstract proof never divides by length or requires a strict
size inequality for successive snapshots.
If the initial entry is absent and every attempt fails,
the final cut remains absent. If the initial entry contains
even a minimal complete snapshot and every attempt fails,
the final cut contains exactly those initial bytes.
If a valid initial entry has a larger byte length than the
first replacement, the reader still sees complete replacement
bytes with no old suffix; the selected object is the staged
object, not a partially overwritten prefix of the old object.
Reversing the lengths gives the symmetric outcome by object
selection, with no old prefix constraint.
At the trace boundary $A_k=0$, $L_k=P_0$ is the defined
fallback. At the other bound $A_k=k$, the abstract count
requires every indexed event to be a publication. With an
initially empty private map and an explicit preparation event
for each commit, this can occur only at cut zero. The invariant
does not assert achievability of its coarse upper bound.
At $T=1$ a nonpublishing event preserves $P_0$; a publishing
event needs a valid closed stage at cut zero, for instance
when the described trace is a suffix of an earlier execution.
An initially present but malformed final file is excluded by
H\ref{hyp:init}. Before a new successful publication, the model
with that relaxed premise can show preservation of the malformed
bytes, but cannot claim eligibility or successful rehydration.


\item \textbf{Multiple writes and serialization limits.}
Order the nonempty commit indices as $i_1<\cdots<i_m$,
where $m=A_T$. For any $k\geq i_m$, the largest committed
index in $C_k$ is $i_m$ and $P_k$ equals the payload there.
For a cut between $i_j$ and $i_{j+1}$, with $1\leq j<m$,
the maximum is $i_j$; hence $P_k$ equals that payload.
When $m=1$ this interior index set is empty, so there is
only the before-first and after-first partition.
When $m=0$ there is no maximum to evaluate,
and the fallback part of Definition~\ref{def:trace} applies.
These cases account for every cut, by comparing $k$
to a finite increasing list of indices.
Physical commit order does not imply chronological order of
snapshot data. If payload $b_1$ has a later application timestamp
than $b_2$, then publishing $b_2$ after $b_1$ still makes
$b_2$ the final contents. The theorem compares transition
indices, not application timestamps.
The process-local lock makes the no-other-attempt clause used
above applicable to cooperating writes within that process.
Two distinct processes can each stage valid bytes in private
temporary directories, and the filesystem may serialize their
namespace replacement events.
Even then no winner is determined by the logical age of the
data in those bytes. For a concrete instance, let both
processes start from the same run version and record different
stage completions, each omitting the other's update.
If they both publish, the later replacement can hide the other
completed stage in subsequent recovery.
Each file can still be whole at every lookup.
If one process fails while the other publishes, the unchanged
public value cannot be attributed to the whole failing time
interval, only to the failure event itself.
The theorem's interval claim therefore uses serialized writers;
a process-local lock supplies no interprocess serialization.
Establishing application-level lost-update freedom would need
version conflict checks or a shared transaction protocol.

\item \textbf{Refinement obligations for the application writer.}
The source-to-model mapping has a specific dependency order.
The call must serialize a complete run snapshot before publishing.
Creating the destination parent directory creates no final
file; any early directory creation is a namespace effect
on a different name from the public entry.
The resulting destination directory must remain stable
through staging and replacement.
A per-attempt private temporary directory inside that parent
supplies the fresh key and normally the same filesystem.
Under the fixed namespace assumption, the destination parent
does not change while the staged file is being written.
Moving the staging directory elsewhere may invalidate
the same-filesystem premise and produce a failed replacement.
The writer invokes the complete staged write, then closes the
staged file before trying to replace the public entry.
A successful write alone does not prove publication,
and a failure in either step cannot reach the replacement.
Publication acts on the staged file, rather than its container.
The cleanup then acts on the temporary directory, rather than
on the public entry.
The literal \texttt{run.json} lookup excludes a temporary
directory and its internal file by construction.
No fallback search for other JSON files is licensed by
H\ref{hyp:read}; adding one would widen the observable namespace.
A logging handler which performs unrelated public file effects
would be an external writer outside the assumed case.
Application correctness additionally needs the values entering
serialization to satisfy the snapshot schema.
For example, a payload that decodes as a JSON number rather
than a run object must not be assigned to $\Valid$ simply
because it has syntactically valid JSON.
The application produces a complete snapshot from its current
run state. The theorem does not make validation guarantees
about unrelated preexisting files or user-added run directories.
Later transformation of an old \texttt{running} status into
an interrupted error is an application rule after decoding.
It does not change which byte string the read selected.

\item \textbf{Independence audit.}
The invariant was derived from finite induction, the effect
of successful replacement, and isolation of other write sets.
There is no assumption about model output quality, the workflow
graph's stage count, a model provider, network availability,
test coverage, or the semantics of SSE delivery.
Those components do not appear in $P_k$, $S_k$, or $C_k$.
There is no appeal to process scheduling fairness or eventually
settled storage.
The writer-serialization assumption is needed for the
no-other-attempt argument over an entire attempt's interval.
It is not needed for the per-transition equality if a given
multiwriter trace independently has ordered atomic namespace
effects and separate staging keys.
The reader immutability premise is needed when the selected
object is read after a replacement; it is not needed to
establish $P_k=L_k$ as a namespace state equation.
Payload validity is used for the membership conclusion,
while the recurrence would still hold for arbitrary bytes.
Private cleanup success is not used in the induction,
which is why the conclusion is about visibility and not
space reclamation.
Acknowledgement is not used to count a commit,
which is why a missing acknowledgement proves neither
that the process published nor that it failed before publication.
A later committed snapshot can have the same application
timestamp or exactly equal bytes without changing the proof.
None of these nondependencies asserts that the corresponding
operational concern is unimportant; it identifies which
additional theorem would be needed to guarantee it.

\item \textbf{Auditor checklist.}
An auditor can check the executable refinement in the
following finite order.
\begin{enumerate}
\item Locate the producer of serialized bytes and establish that
it completes before the replacement. Check the explicit encoding.
\item Trace every write and cleanup target. Resolve the intended
public path and verify that every other path is private to
one attempt under namespace assumptions.
\item Confirm that write and close failures cannot reach the
replacement. Exercise both absent and existing final names.
\item Confirm that the successful replacement is the only
operation that rebinds the public entry. Check that the staged
file is in the destination filesystem and closed before use.
\item Inject a failure before replacement and compare the old
and observed bytes exactly. Inject a cleanup failure separately
and identify whether it follows publication.
\item Exercise recovery against an orphaned stage without a
public file, and beside an older valid public file. Verify the
loader's literal name projection and rehydration rule separately.
\item Check the running environment's namespace guarantees.
An operating system may reject replacement while a reader holds
an incompatible open handle. That rejection warrants a failure
case, not a fictional successful concurrent publication.
\item Record the scope of the observations. Test results on one
local environment do not independently validate all platform
or mount choices.
\end{enumerate}

\item \textbf{Scope of the established result.}
The observation equation is a visibility statement for one
file and successful selections.
It gives no upper bound on I/O latency and no guarantee that
a replacement will ever succeed.
Restart recovery after process interruption is distinguished
from recovery after power loss, filesystem corruption, remote
storage anomalies, or external manipulation of the directory.
The proof contains no file or directory \texttt{fsync} and
does not establish flushing order or crash durability.
It also does not combine the sidecar and specialist artifacts
into a transaction. A reader may recover the last sidecar
snapshot without all related stage files matching that moment.
The proof supplies neither a cross-run snapshot nor rollback
of a stage or publication authorization.
Within the stated scope, the induction and selection arguments
establish all parts of Theorem~\ref{thm:visibility}.
\end{enumerate}
\end{proof}

\section{Why direct truncation is insufficient}
\label{sec:countermodel}
\begin{counterexample}[A complete old snapshot can be lost]
\label{cex:truncate}
Let $P_0=b_0$ be a complete application snapshot.
A writer opens the final name in truncating mode.
After truncation and before a complete write, the final entry
can contain $\epsilon$ or a proper prefix $p$ of the desired
payload. In particular choose a cut immediately after truncation
and interrupt the writer there. The final entry now holds
$\epsilon\notin\Valid$ even though no complete publication
occurred. Consequently a recovery lookup at that cut cannot
recover $b_0$. A lock that serializes writers leaves this
counterexample intact, since it has only one writer.
\end{counterexample}

\section{Local verification boundary}
\label{sec:verification}
Compare the pre-failure and post-failure bytes in the same run
directory. Exercise failures in the staged write and replacement,
and verify repeated publication and literal-name recovery.
Successful local tests constrain the implementation observed in
those tests. The independent reviewer should compare the source,
test receipts, and compiled paper to the final candidate.
File and directory flush behavior, cross-platform durability,
and the integrity of remote mounts require evidence outside this
companion.

\begin{thebibliography}{9}
\bibitem{python}
Python Software Foundation, \textit{Python Standard Library},
\texttt{os.replace} and \texttt{tempfile.TemporaryDirectory},
Python 3 documentation.

\bibitem{posix}
The Open Group, \textit{POSIX.1-2024, System Interfaces},
\texttt{rename}, including replacement visibility of the
destination name within the filesystem assumptions.
\end{thebibliography}

\clearpage
\section*{Changelog}
\addcontentsline{toc}{section}{Changelog}
\begin{table}[!htbp]
\centering\small
\caption{Revision history}
\label{tab:changelog}
\begin{tabularx}{\linewidth}{@{}l l X@{}}
\toprule
Version & Date & Change \\
\midrule
2026.1.0.0 & 2026-11-20 & Introduced the single-run namespace model,
trace proof, direct-write counterexample, and implementation audit. \\
\bottomrule
\end{tabularx}
\end{table}
\end{document}
